%\pagestyle{mystyle}
\chapter{Knowledge Extraction in the Legal Domain}
\todo{kg because in wise it is not only entities but also relations from chat metadata}
\todo{add dave}
% \section{User interfaces}
% \subsection{Document Explorer}
% \subsection{Knowledge Editing to Correct Machine Errors}
% Human-in-the-Loop
% \subsection{Faceted Search}
% \subsection{Question Answering with Retrieval Augmented Generation}

% \section{Question Answering with ReFactX}

As anticipated in \Cref{sec:background:ai4legal}, \ac{ai} systems can be applied to the legal domain for addressing the analysis of vast amount of documents, that can come from heterogeneous sources.
During my Ph.D. I worked with two types of documents from the Italian legal domain: civil judgements and investigative chat logs. In both cases, the main problem addressed is the difficulty that the users---such as judges or prosecutors---face in analyzing the high number of records. Indeed, my co-authored contributions seek to empower final users with tools that allow them to efficiently explore and retrieve relevant documents or chat messages.
These proposals have been published in Computer Law \& Security Review (Vol. 52)~\cite{BELLANDI2024}, in the proceedings of the 22nd International Conference of the Italian Association for Artificial Intelligence (AIxIA 2023 -- Advances in Artificial Intelligence)~\cite{aixia_2023}, and of the 25th International Conference on Web Information Systems Engineering (WISE 2024)~\cite{pozzi_wise_2025}.

Data in this domain, as described in \Cref{sec:background:ai4legal}, often contain personal data, and privacy regulations implies important challenges such as the difficulty in obtaining labeled data. Furthermore, the language used in documents, such as civil judgments, differs from the one generally used for training neural models---i.e., Web data~\cite{}.\todo{giustifica. guarda miei paper, legal bert,...}

\todo{private data challenge and consequences: no api,...}
\todo{summarize abstracts and focus on ie (ia later)}
\todo{contributions itemize}

Labeling data is a problem\todo{so gold standard annotation}
Highlight the challenges (e.g. data,...)
\todo{objective and motivation of IE} enable ia techniques (faceted search, sta/browsing annotated docs)
semantic search, sta, kbp (more wise)
\todo{use background to motivate. start from general IA for law. then IE. then IE contributions. finally IA contributions }
\begin{itemize}
    \item context
    \item challenges of legal domain
    \item ie motivation (sta, for ia). vast documents
    \item ie contributions
    \item ia intro
    \item ia contributions
\end{itemize}